\chapter*{\sffamily 2024年\textsf{\textsc{Hdu}}「概率论 \& 数理统计」期中试题}

\ExplSyntaxOn
\vspace{-12.0pt}\setcounter{page}{1}
\begin{center}
  \begin{minipage}{.66\textwidth}
    \centering
    \begin{tasks}
      [ label-width = 5\ccwd, item-indent = 2cm,
        after-item-skip = 1ex ](2)
      \task[\textbf{课程信息}：] 理学院/A0714030
      \task[\textbf{解析提供}：] 幻想大白
      \task[\textbf{授课教师}：] 大学数学教学团队
      \task[{\makebox[2em][c]{\textbf{\textsf\TeX}}}\textbf{排版}：] \textsf{Izumi~Sakai}
    \end{tasks}
  \end{minipage}
  \begin{minipage}{.15\textwidth}
    \centering
    \begin{tikzpicture}
      \node [ inner~sep = 0 ] (wechat) {
        \fancyqr
          [ size = 4em, image~padding = 1,
            image = \raisebox{4em/33}{\textcolor{DarkGreen \l__whats_darkmode_fading_tl}{\faIcon{weixin}}},
            r~color = DarkGreen \l__whats_darkmode_fading_tl, version = 4
          ]{http://weixin.qq.com/r/sRJiemTEGjllrSAZ90ej}
        };
      \node [ inner~sep = .5ex, align = center, right, font = \scriptsize ]
        at (wechat.east) {未\\央\\学\\社};
    \end{tikzpicture}
  \end{minipage}
  \begin{minipage}{.15\textwidth}
    \centering
    \begin{tikzpicture}
      \node [ inner~sep = 0 ] (qq) {
        \fancyqr
          [ size = 4em, image~x~padding = 2, image~y~padding = 1,
            image = \raisebox{4em/33}{\textcolor{MidnightBlue \l__whats_darkmode_fading_tl}{\faIcon{qq}}},
            l~color = MidnightBlue \l__whats_darkmode_fading_tl, version = 4
          ]{https://qm.qq.com/q/rbfYuqPhQs}
        };
      \node [ inner~sep = .5ex, align = center, left, font = \scriptsize ]
        at (qq.west) {数\\学\\营};
    \end{tikzpicture}
  \end{minipage}
\end{center}
\ExplSyntaxOff

\section{选择题（每题3分，共30分）}

\begin{problem}
  以 $A$ 表示``甲种产品合格，乙种产品合格''，则其对立事件 $\overline A$ 为
  \begin{tasks}(2)
    \task 甲种产品不合格，乙种产品合格
    \task 甲、乙两种产品均不合格
    \task 甲种产品合格、乙种产品不合格
    \task[\true] 甲种产品不合格或乙种产品不合格
  \end{tasks}
\end{problem}
\begin{solution}
  设甲、乙产品合格为 $A_1,~ A_2$. $A = A_1 \cup A_2$.
  $\overline A = \overline{A_1 \cup A_2} = \overline A_1 \cup \overline A_2$.
  所以D选项正确.
\end{solution}

\begin{problem}
  设 $A,~ B,~ C$ 为三个随机事件，则 $A,~ B,~ C$ 至少有一个不发生的概率为
  \begin{tasks}(4)
    \task $\overline A\, \overline B\, \overline C$
    \task[\true] $\overline{ABC}$
    \task $A \cup B \cup C$
    \task $ABC$
  \end{tasks}
\end{problem}
\begin{solution}
  $A$、$B$、$C$ 至少一个不发生为
  $\overline A \cup \overline B \cup \overline C = \overline{ABC}$.
\end{solution}

\begin{problem}
  设 $A$、$B$ 为相互独立的随机事件，$P(\overline AB) = 3P(A)P(B) > 0$，则 $P(A|B) = $
  \begin{tasks}(4)
    \task $1$
    \task $1/2$
    \task $1/3$
    \task[\true] $1/4$
  \end{tasks}
\end{problem}
\begin{solution}
  $P(\overline A B) = P(B) - P(A)P(B) = 3P(A)P(B)$，所以
  $P(B) = 4P(AB),~ P(A|B) = P(AB)/P(B) = 1/4$.
\end{solution}

\begin{problem}
  设 $F_1(x),~ F_2(x)$ 分别为随机变量 $X_1$ 和 $X_2$ 的分布函数.
  为使 $aF_1(x) - bF_2(x)$ 为某一随机变量的分布函数，则在下列给出的各组数值中应取
  \begin{tasks}(4)
    \task $a = 1/3,~ b = 1/3$
    \task $a = -1/2,~ b = 3/2$
    \task[\true] $a = 2/5.~ b = -3/5$
    \task $a = 1/2.~ b = -3/2$
  \end{tasks}
\end{problem}
\begin{solution}
若 $F_1(x),~ F_2(x),~ \cdots,~ F_n(x)$ 都是分布函数，
当 $k_i \geq 0$ 且 $\sum k_i = 1$ 时，$\sum_{k = 1}^n k_i F_i(x)$ 也是分布函数.
\end{solution}

\begin{problem}
  设连续型随机变量 $X$ 的概率密度函数和分布函数分别为 $f(x)$ 和 $F(x)$.
  则根据 $f(x)$ 和 $F(x)$ 的性质，下列说法正确的是
  \begin{tasks}(4)
    \task[\true] $f(x)$ 可以是偶函数
    \task $f(x)$ 可以是奇函数
    \task $F(x)$ 可以是偶函数
    \task $F(x)$ 可以是奇函数
  \end{tasks}
\end{problem}
\begin{solution}
  由 $f(x)$ 的性质 $\int_{-\infty}^{+\infty} f(x) \d x = 1$，
  若 $f(x)$ 为奇函数，则 $\int_{-\infty}^{+\infty} f(x) \d x = 0$，所以B选项错误.
  由 $F(x)$ 的性质
  $\lim_{x \to -\infty} F(x) = 0,~ \lim_{x \to +\infty} F(x) = 1$，
  $F(x)$ 不可能为奇函数或偶函数.
\end{solution}

\begin{problem}
  设随机变量 $X$ 和 $Y$ 相互独立，分别服从参数为 $-\frac12$ 和 $\frac12$ 的指数分布
  （指数分布密度函数 $F(x) =
  \begin{cases}
    \frac1\theta e^{-x/\theta}, & x > 0\\
    0,                          & x \leq 0
  \end{cases}$，$\theta$ 为参数），则 $P\{X > Y\} = $
  \begin{tasks}(4)
    \task $1/2$
    \task $1/6$
    \task $2/5$
    \task[\true] $3/5$
  \end{tasks}
\end{problem}
\begin{solution}
  参考 Problem 19. $P\{X > Y\} = 3/(2 + 3) = 3/5$.
\end{solution}

\begin{problem}
  设随机变量 $X$ 的概率密度为 $f_X(x),~ x \in \mathbb R$. 令 $Y = 1 - 2X$，
  则 $Y$ 的概率密度 $f_Y(y)$ 为
  \begin{tasks}(4)
    \task $-2f_X\ab(1 - \frac y2)$
    \task[\true] $\frac12 f_X\ab(\frac{1-y}{2})$
    \task $2f_X(1 - 2y)$
    \task $-\frac12 f_X\ab(-\frac{1 - y}{2})$
  \end{tasks}
\end{problem}
\begin{solution}
  $  F_Y(y) = P(1 - 2X \leq y) = P\ab(X \geq \frac12(1 - y))
   = \int_{(1 - y)/2}^{+\infty} f_X(x) \d x$.
  $f_Y(y) = F_Y'(y) = \frac12 f_X\ab(\frac{1 - y}{2})$.
\end{solution}

\begin{problem}
  设随机变量 $X$ 服从正态分布 $N \sim (\mu, \sigma^2)$，则随着 $\sigma > 0$ 减少，
  概率 $P\{|X - \mu| > \sigma\}$
  \begin{tasks}(4)
    \task 单调减少
    \task 单调增加
    \task[\true] 保持不变
    \task 减少或增加不确定
  \end{tasks}
\end{problem}
\begin{solution}
  $P(|x - \mu| > \sigma) = P\ab(\frac{|x - \mu|}{\sigma} > 1) = 1 - \Phi(1)$，
  所以概率 $P\{|X - \mu| > \sigma\}$ 保持不变.
\end{solution}

\begin{problem}
  设随机变量 $X$ 和 $Y$ 相互独立，服从相同的分布，且 $X$ 的分布函数为 $F(x)$，
  则 $Z = \max(X, Y)$ 的分布函数为
  \begin{tasks}(4)
    \task[\true] $F^2(x)$
    \task $F(x)F(y)$
    \task[\true] $F^2(z)$
    \task $1 - [1 - F(z)]^2$
  \end{tasks}
\end{problem}
\begin{solution}
  $ F_Z(z) = P\{Z \leq z\} = P{\max\{X, Y\} \leq z} = P\{X \leq z,~ Y\leq z\}
  = P\{X \leq z\} P\{Y \leq z\} = F(x)F(y) = F^2(x)$.
\end{solution}

\begin{problem}
  设随机变量 $(X, Y)$ 在单位元内服从均匀分布，其概率密度函数为 $f(x,y) =
  \begin{cases*}
    C, & $x^2 + y^2 < 1$\\
    0, & else
  \end{cases*}$，则 $C = $
  \begin{tasks}(4)
    \task $1$
    \task $2$
    \task $\pi$
    \task[\true] $1/\pi$
  \end{tasks}
\end{problem}
\begin{solution}
  二维随机变量在区域 $D:~ x^2 + y^2 < 1$ 服从均匀分布，所以 $C = 1/S = 1/\pi$.
\end{solution}

\section{填空题（每题3分，共15分）}

\begin{problem}
  甲口袋有 $5$ 个白球、$3$ 个黑球，乙口袋有 $4$ 个白球、$6$ 个黑球，
  从两口袋中任取一球，取到的两个球颜色相同的概率为 \result{$19/40$}.
\end{problem}
\begin{solution}
  $P = \frac58 \times \frac4{10} + \frac38 \times \frac6{10} = \frac{19}{40}$.
\end{solution}

\begin{problem}
  设 $S = \{x | 0 \leq x \leq 2\},~ A = \ab\{A | \frac12 < x \leq 1\},~
  B = \ab\{x | \frac14 \leq x < \frac23\}$，
  则 $\overline{\overline A\,\overline B} = $
  \result{$\{x | \frac14 \leq x \leq 1\}$}.
\end{problem}
\begin{solution}
  $\overline{\overline A\,\overline B} = A \cup B = \{x|1/4 \leq x \leq 1\}$.
\end{solution}

\begin{problem}
  设随机变量 $X$ 的分布律为 $P\{X = k = A(2 + k)^{-1}\},~ k = 0,~ 1,~ 2$，
  则常数 $A = $ \result{$12/13$}.
\end{problem}
\begin{solution}
  $ \sum_{k = 0}^2 P\{X = k\} \frac12 k + \frac13 k + \frac14 k = 1,~
    k = \frac{12}{13}$.
\end{solution}

\begin{problem}
  设随机变量 $X$ 的概率密度函数为 $f(x) = \frac{1}{3\sqrt{2\pi}}e^{-(x + 1)^2/18},~
  x \in \mathbb R$，则 $Y = $ \result{$(x + 1)/3$} $\sim N(0,1)$.
\end{problem}
\begin{solution}
  由正态分布概率密度函数
  $f(x) = \frac{1}{\sqrt{2\pi}\sigma}e^{-(x - \mu)^2/2\sigma^2}$
  得 $\mu = -1,~ \sigma = 3$.
  所以 $Y = (x - \mu)/\sigma = (x + 1)/3 \sim N(0,1)$.
\end{solution}

\begin{problem}
  设 $(X,Y)$ 为二维连续型随机变量，则 $P\{X = Y\} = $ \result{$0$}.
\end{problem}
\begin{solution}
  二维连续型随机变量落在它平面区域内任意一条直线或曲线上概率均为 $0$.
\end{solution}

\section{计算题（共49分）}

\begin{problem}[本题5分]
  根据美国的一份资料显示，在美国总的来说患肺癌的概率约为 $0.1\%$，
  在人群中有 $20\%$ 是吸烟者，他们患肺癌的概率约为 $0.4\%$，
  利用全概率公式求不吸烟者患肺癌的概率是多少？
\end{problem}
\begin{solution}
  令``患肺癌事件''为 $A$，$B_1,~ B_2$ 分别表示吸烟者和不吸烟者.
  则 $P(A) = P(A|B_1)P(B_1) + P(A|B_2)P(B_2)$.
  又 $P(A) = 0.1\%,~ P(A|B_1) = 0.4\%,~ P(B_1) = 20\%,~ P(B_2) = 80\%$.
  代入得 $P(A|B) = 0.025\%$.
\end{solution}

\begin{problem}[本题10分]
  设随机变量 $X$ 的分布函数为 $F(x) =
  \begin{cases}
    0,    & x < 0\\
    ax^2, & 0 \leq x \leq 1\\
    1,    & x > 1
  \end{cases}$.
  \begin{tasks}[ label = \arabic*. ](3)
    \task 求 $a$ 的值.
    \task 求 $X$ 的密度函数 $f(x)$.
    \task 求 $P\{0.3 < X < 0.7\}$.
  \end{tasks}
\end{problem}
\begin{solution}\leavevmode
  \begin{tasks}[ label = \arabic*. ](2)
    \task $\lim_{x \to 1^+} F(x) = F(1)$，即 $1 = F(1) = ax^2|_{x = 1} = a$.
    \task $f(x) = F'(x) =
    \begin{cases}
      2x, & 0 \leq x \leq 1\\
      0,  & \text{else}
    \end{cases}$.
    \task! $P\{0,~ 3 < x < 0\} = F(0.7) - F(0.3) = 0.7^2 - 0.3^2 = 0.4$.
  \end{tasks}
\end{solution}

\begin{problem}[本题12分]
  设 $X,~ Y$ 是相互独立的随机变量，其密度函数分别为
  \[
    f_X(x) =
    \begin{cases}
      \alpha e^{-\alpha x}, & x > 0\\
      0,                    & x \leq 0
    \end{cases},~
    f_Y(x) =
    \begin{cases}
      \beta e^{-\beta x}, & y > 0\\
      0,                  & y \leq 0
    \end{cases}
  \]
  其中常数 $\alpha > 0,~ \beta > 0$. 定义随机变量
  $Z =
  \begin{cases}
    1, & \text{当 $X \leq Y$}\\
    0, & \text{当 $ X > Y$}
  \end{cases}$.
  \begin{tasks}[ label = \arabic*. ](3)
    \task 求条件概率 $f_{X|Y}(x|y)$.
    \task 求 $Z$ 的分布律.
    \task 求 $Z$ 的分布函数.
  \end{tasks}
\end{problem}
\begin{solution}\leavevmode
  \begin{enumerate}
    \item 由于 $X,~ Y$ 独立，所当 $y > 0$ 时，
    $f_{X|Y}(x,y) = f_X(x) =
    \begin{cases}
      \alpha e^{-\alpha x}, & x > 0\\
      0,                    & x \leq 0
    \end{cases}$.
    \item $f(x,y) = f_X(x) \cdot f_Y(y) =
    \begin{cases}
      \alpha \beta e^{-\alpha x} e^{-\beta y}, & x > 0,~ y > 0\\
      0,                                       & \text{else}
    \end{cases}$. 积分区域 $D:\{0 < x < +\infty,~ 0 \leq x \leq y\}$.
    \begin{align*}
    & P\{X > Y\}
    = \iint_D f(x,y)\d x \d y
    = \int_0^{+\infty} \alpha e^{-\alpha x}\d x \int_0^x \beta e^{-\beta y}\d y
    = \frac{\beta}{\alpha + \beta}\\
    & P\{X < Y\} = 1 - P\{X > Y\} = \frac{\alpha}{\alpha + \beta}
    \end{align*}
    \begin{minipage}{.54\linewidth}
      \centering
      \begin{minipage}{.36\linewidth}
        \centering
        \begin{tikzpicture} 
          \draw [ -> ] (-.2,0) -- (1.8,0) node [ below ] {$x$};
          \draw [ -> ] (0,-.2) -- (0,1.8) node [ left ] {$y$};
          \fill [ pattern = north west lines ]
            (0,0) -- (1.5,0) -- (1.5,1.5) -- cycle;
          \draw (-.2,-.2) -- (1.5,1.5) node [ left ] {$y = x$};
        \end{tikzpicture}
      \end{minipage}
      \hfill
      \begin{minipage}{.6\linewidth}
        \centering
        \begin{tabular}{*3c}
          \toprule
          $Z$ & $0$ & $1$\\
          \midrule
          $P$ & $\dfrac{\beta}{\alpha + \beta}$
              & $\dfrac{\alpha}{\alpha + \beta}$\\
          \toprule
        \end{tabular}
      \end{minipage}
    \end{minipage}
    \begin{minipage}{.42\linewidth}
      \item $F_Z(z) = P\{Z \leq z\} =
      \begin{dcases}
        0,                            & z < 0\\
        \frac{\beta}{\alpha + \beta}, & 0 \leq z < 1\\
        1,                            & z \geq 1
      \end{dcases}$.
    \end{minipage}
  \end{enumerate}
\end{solution}

\begin{problem}[本题12分]
  抛一枚硬币 $3$ 次，$X,~ Y$ 分别表示出现 $H$ 的次数和出现 $H,~ T$ 次数差的绝对值.
  \begin{tasks}[ label = \arabic*. ](7)
    \task*(2) 求 $(X,Y)$ 的分布律.
    \task*(2) 求分布函数 $F(2.5,1.5)$ 值.
    \task*(3) 问 $X$ 和 $Y$ 是否相互独立，并说明理由.
  \end{tasks}
\end{problem}
\begin{solution}\leavevmode\newline
  \begin{textfig}[5:8]\leavevmode
    \begin{enumerate}
      \renewcommand*\arraystretch{.5}
      \item 
      \begin{tabular}{*5c}
        \toprule
        $Y \backslash X$ & $0$   & $1$   & $2$   & $3$ \\
        \midrule
        $1$              & $0$   & $3/8$ & $3/8$ & $0$\\
        $3$              & $1/8$ & $0$   & $0$   & $1/8$\\
        \bottomrule
      \end{tabular}
    \end{enumerate}
    \textfigsplit
    \begin{enumerate}
      \addtocounter{enumi}{1}
      \item $F(2.5,1.5) = P\{X = 1,~ Y = 1\} + P\{X = 2,~ Y = 1\} = 3/4$.
      \item $P\{X = 0,~ Y = 1\} \neq P \{X = 0\}P\{Y = 1\}$，
      $X,~ Y$ 不相互独立.
    \end{enumerate}
  \end{textfig}
\end{solution}

\begin{textfig}[3:1]
  \begin{problem}[本题10分]
    设随机变量 $X$ 与 $Y$ 的联合分布律如右侧图表所示.
    \begin{tasks}[ label = \arabic*. ](2)
      \task 求 $p$.
      \task 求 $P\{X < 1 | Y = 2\}$.
      \task 求 $P\{X = 2Y\}$.
      \task 求 $Z = X + Y $ 的分布函数.
    \end{tasks}
  \end{problem}
  \textfigsplit\centering
  \renewcommand*\arraystretch{.6}
  \begin{tabular}{cccc}
    \toprule
    $Y \backslash X$ & $0$   & $1$ & $2$\\
    \midrule
    $0$                    & $1/4$ & $0$ & $1/4$\\
    $1$                    & $0$   & $p$ & $0$\\
    $2$                    & $1/12$& $0$ & $1/12$\\
    \bottomrule
  \end{tabular}
\end{textfig}
\begin{solution}\leavevmode\newline
  \begin{textfig}[1:1]\leavevmode
    \begin{enumerate}
      \item $1/4 + 1/4 + p + 1/12 + 1/12 = 1,~ p = 1/3$.
      \item $P\{X < 1|Y\} = P\{X < 1,~ Y = 2\}/P\{Y = 2\} = 1/2$.
      \item $P = P\{X = 0,~ Y = 0\} + P\{X = 2,~ Y = 1\} = 1/4$.
    \end{enumerate}
    \textfigsplit\centering
    \begin{enumerate}
      \addtocounter{enumi}{3}
      \item 若保留 $X + Y = 1,~ 3$，扣1分.
      
      \renewcommand*\arraystretch{.8}
      \begin{tabular}{cc}
        \toprule
        $Z$ & $P$\\
        \midrule
        $0$ & $1/4$\\
        $2$ & $2/3$\\
        $4$ & $1/12$\\
        \bottomrule
      \end{tabular}
      $F_{X + Y}(z) =
      \begin{cases}
        0,     & z < 0\\
        1/4,   & 0 \leq z < 2\\
        11/12, & 2 \leq z < 4\\
        1,     & z \geq 4
      \end{cases}$.
    \end{enumerate}
  \end{textfig}
\end{solution}

\begin{problem}[本题6分]
  设随机变量 $X,~ Y$ 相互独立，且 $X \sim \pi(\lambda_1),~ Y \sim \pi(\lambda_2)$.
  证明：$X + Y \sim \pi(\lambda_1 + \lambda_2)$.
\end{problem}
\begin{solution}
  \begin{proof}
    $X \sim \pi(\lambda_1),~ Y \sim \pi(\lambda_2)$ 相互独立，
    $ P\{X = k\} = \frac{e^{-\lambda_1} \lambda_1^k}{k!},~
      P\{Y = l\} = \frac{e^{-\lambda_2} \lambda_2^k}{l!},~ k,~ l = 0,~ 1,~ 2,~ \cdots$.
    $X + Y$ 所有能取的值为 $n = 0,~ 1,~ 2,~ 3,~ \cdots$. 所以
    \begin{align*}
        & P\{X + Y = n\}
      = \sum_{k = 0}^n P\{X = k\} P\{Y = n - k\}
      = \sum_{k = 0}^n \frac{e^{-\lambda_1} \lambda_1^k}{k!}
                       \frac{e^{-\lambda_2} \lambda_2^k}{l!}\\
      = & \frac{e^{-(\lambda_1 + \lambda_2)}}{n!}
        \sum_{k = 0}^n \frac{n!}{k!(n - k)!} \lambda_1^k \lambda_2^{n - k}
      = \frac{e^{-(\lambda_1 + \lambda_2)}}{n!}
        \sum_{k = 0}^n C_n^k \lambda_1^k \lambda_2^{n - k}
      = \frac{e^{-(\lambda_1 + \lambda_2) \cdot (\lambda_1 + \lambda_2)^n}}{n!}
    \end{align*}
    即 $X + Y \sim \pi(\lambda_1 + \lambda_2)$.
  \end{proof}
\end{solution}